\documentclass{article}
\usepackage{neurips_2026}
\usepackage[utf8]{inputenc}
\usepackage[T1]{fontenc}
\usepackage[scaled=0.9]{couriers}
\usepackage{hyperref}
\usepackage{url}
\usepackage{booktabs}
\usepackage{amsfonts}
\usepackage{amsmath}
\usepackage{nicefrac}
\usepackage{microtype}
\usepackage{xcolor}
\usepackage{graphicx}

\title{Minimum Targeted Investment to Top a Multi-Task Leaderboard Depends on the Aggregation Rule}

\author{Anonymous Author(s)}

\begin{document}
\maketitle

\begin{abstract}
Public leaderboards rank models by combining scores on many tasks, and developers can raise a model's scores on chosen tasks.
We ask how much targeted improvement a model needs to reach first place and how this depends on the aggregation rule.
We formalise a score counterfactual in which one target raises its scores up to a cap while all rivals stay fixed, and report three separate optima (total points, tasks touched, largest single-task lift) and, separately, whether first place is attainable at all.
We give closed forms for the mean, median and minimum, and solve the geometric mean exactly and the Borda and Condorcet-winner requirements by integer programming.
On the Open LLM Leaderboard (4,495 targets, six tasks) and HELM MMLU (78 targets, 57 subjects), first place is out of reach for most models if no task may be raised by more than 10 points: it is attainable for 3--29\% of Open LLM Leaderboard targets and 18--51\% of HELM targets, depending on the rule.
The rules differ in how many models can win, not only in cost; on the cohort feasible under every rule with loose caps, the minimum rule needs the fewest points and Borda the most tasks.
The ordering of rules changes between datasets and caps, so our conclusions are conditional on the cap and the score scale.
Investment here is a score counterfactual, not a training cost.
\end{abstract}

\section{Introduction}
\label{sec:intro}

Public leaderboards turn a score matrix of models by tasks into a single ranking, and the top place carries visibility.
The ranking depends on the aggregation rule, and developers can raise a model's score on chosen tasks, by targeted data or tuning.
We ask a counterfactual question about the scores: \emph{how much targeted improvement does a given model need to reach first place, and how does that depend on the rule that aggregates the tasks?}
We answer it for a single target model with all rivals held fixed, on two public matrices, the Open LLM Leaderboard (OLLB; 4{,}496 models, six tasks) and HELM MMLU (79 models, 57 subjects).

The question has three parts that earlier work on task selection and aggregation does not separate (Section~\ref{sec:related}).
First, ``how much'' is not one number: we report the total points added ($L_1^\ast$), the number of tasks touched ($K^\ast$) and the largest single-task lift ($L_\infty^\ast$) as three separate optima.
Second, an investment is meaningful only under a cap on how far each score can rise, and under attainable caps first place is often out of reach; we therefore report \emph{feasibility} (whether first place is attainable at all, and the smallest cap or uniform lift that makes it so) as the headline, and compare costs only on the cohort of targets that are feasible under every compared rule.
Third, the answer may depend on the score scale and the number of tasks, which we vary.

\paragraph{Contributions.}
(i) We formalise the score-edit counterfactual and its three optima, with feasibility kept separate from cost, for score rules (mean, median, minimum, geometric mean) and ordinal rules (Borda, strict Condorcet winner; the latter is an upper bound for any Condorcet-consistent rule).
(ii) We give closed forms for the mean, median (odd $n$) and minimum, and show that $L_\infty^\ast$ equals the smallest uniform lift for every rule; we check them against exact solvers, including integer programs for the ordinal rules, on 31{,}080 instances, all of which agree (Section~\ref{sec:method}, Appendix~\ref{app:proofs}).
(iii) We run all rules over eight cap scenarios on 4{,}495 OLLB and 78 HELM targets (Section~\ref{sec:exp}).

\paragraph{Findings.}
They are conditional on the cap and the scale and are stated as such.
\textbf{(F1)} If no task may be raised by more than 10 points, first place is attainable for only 3--29\% of OLLB targets and 18--51\% of HELM targets, depending on the rule; for models ranked 2--30 under the mean the shares are 76--100\% and 41--100\%.
Rules differ in \emph{how many} models can win, not only in what it costs: on OLLB, Borda reaches 29\% at 10 points where the mean reaches 3\%.
\textbf{(F2)} The ordering of rules is not stable across datasets and caps: the minimum rule is the hardest to satisfy under small caps on HELM (18\% at 10 points) and the easiest on OLLB under a 20-point cap (64\% against 14--21\% for the mean, median and geometric mean).
\textbf{(F3)} Under loose caps, on the cohort of targets feasible under every rule, the minimum rule needs the fewest total points on OLLB and, with the median, on HELM, while Borda needs the most tasks (5 of 6 on OLLB, 46 of 57 on HELM).
\textbf{(F4)} On OLLB, whose six tasks have unrelated scales, the number of tasks needed by targets ranked beyond 10 changes by up to four with the score scale (official, raw, $z$-score, rank percentile), so OLLB conclusions are tied to the official scale.
\textbf{(F5)} The closed forms explain these differences: the minimum rule requires every task below the best rival's level to reach it, so it is blocked by any binding cap there.

\emph{Investment here is a score counterfactual, not a training cost}; it ignores spillover between tasks, rival responses and the data or compute that would deliver the gain (Section~\ref{sec:limits}).

\section{Related Work}
\label{sec:related}

\paragraph{Aggregation and task selection in multi-task evaluation.}
Multi-task results depend on which tasks are included and how they are combined.
\citet{dehghani2021lottery} show that rankings can change with task selection and discuss macro-averaging, arithmetic versus geometric means and rank-based aggregation; \citet{liang2023helm} leave a universal single-number aggregate explicitly unresolved and stress that it encodes value judgements.
These works vary the benchmark while the compared models stay as they are. We instead hold the benchmark and all rivals fixed and edit one model's scores.
They do not establish an ordering of rules by the cost of changing a ranking, and we do not read one into them.

\paragraph{Social choice and evaluation.}
\citet{lanctot2023vase} cast evaluation as voting, with models as alternatives and tasks as voters, so that comparisons need only within-task order. The Handbook of Computational Social Choice \citep{brandt2016handbook} supplies the voting-rule and Condorcet vocabulary we use \citep{zwicker2016voting}.
Our target-promotion problem is related to priced bribery, where a fixed set of voters is paid to change their votes \citep{faliszewski2016control}; it is a structured special case in which the target may only move up on the tasks it is allowed to edit, so complexity results for unrestricted bribery do not transfer automatically.
We make no claim of new algorithms for bribery or control.

\paragraph{Leaderboard incentives and strategic post-training.}
\citet{singh2025illusion} document private testing, selective disclosure and data-access asymmetries on Chatbot Arena, and find in a controlled training-mixture experiment that gains on the targeted evaluation need not transfer broadly.
\citet{chen2026incentives} analyse simultaneous developers' effort and a designer's tuning policy as an equilibrium problem, under assumptions on the score response to effort.
Our setting is different and narrower: a single target, fixed rivals, observed score matrices, and no equilibrium, dynamics or latent capability.
It measures the \emph{score improvement} that a ranking goal requires under a stated cap, not the training effort that would deliver it.

\paragraph{Benchmark-related training and contamination.}
\citet{zhou2023cheater} show that continued training on benchmark-related data raises benchmark scores, with possible losses elsewhere, and \citet{yang2023rephrased} show that rephrased test samples can evade common decontamination while still inflating scores.
These works motivate the intervention class we study; they do not provide a task-response or cost law, and they show that independent monotone gains with no side effects are an abstraction (Section~\ref{sec:limits}).
Score tables alone neither identify nor exclude contamination, and nothing here accuses any listed model.

\paragraph{Benchmarks and data used.}
We use the HELM MMLU release \citep{liang2023helm,hendrycks2021mmlu} and the Open LLM Leaderboard, whose tasks include BBH \citep{suzgun2022bbh} and which is built on the Evaluation Harness \citep{sutawika2023harness}.

\section{Problem Setting and Method}
\label{sec:method}

\paragraph{Leaderboard and rules.}
A leaderboard is a score matrix $x\in[0,100]^{M\times n}$ of $M$ models on $n$ tasks (sub-datasets).
An aggregation rule $R$ maps a matrix to a ranking of the models.
We study score rules, which compute a scalar per model: the arithmetic \emph{mean} (the equal-task mean; on the Open LLM Leaderboard, the mean of the six official normalised task scores is the official average, while the raw-accuracy columns are a separate variant; we do not claim that it reproduces HELM's own aggregate), the \emph{median}, the \emph{minimum}, and the \emph{geometric mean}; and two ordinal rules, which use only within-task order: the \emph{Borda} count and the requirement to be a strict \emph{Condorcet winner} (beating every rival in pairwise comparison, where a task won counts 1, a task tied counts $\tfrac12$, and the target wins if its total exceeds $n/2$; with ties this is not the same as winning more than $n/2$ tasks strictly).

\paragraph{Intervention and investment.}
Fix a target model $i$. The developer may raise its scores to any $y_t\in[x_{it},c_t]$, where $c_t$ is a cap on task $t$; all rivals' raw scores stay fixed, there are no spillovers to other tasks, and no rival responds.
After the edit, ordinal aggregates of rivals are recomputed because the target may pass them.
The target is \emph{first} if it beats every rival under $R$ by a margin $\varepsilon$ (score rules; $\varepsilon=10^{-3}$ points unless stated) or strictly (ordinal rules).
Strict victory over a continuous score can have an infimum that is not attained, so $\varepsilon$ defines the margin version we report, and we report sensitivity to it; a target that already wins has cost zero.

\emph{Investment is a score counterfactual, not a training cost.}
We measure it by three separate optima, each its own optimisation problem: $L_1^\ast$, the least total number of points added; $K^\ast$, the fewest tasks touched; and $L_\infty^\ast$, the smallest largest single-task lift when any subset may be raised.
Where a solver minimises one of them, the others are only reported descriptors at its solution.
Raising all tasks by a common amount is a different quantity and is not called $L_\infty^\ast$.
A target is \emph{infeasible} if no allowed target row under the cap satisfies the winning condition (a cap that binds at the optimum does not make a target infeasible); we report the feasible fraction separately from costs over feasible targets, never averaging an infinite cost, together with the smallest cap at which the target becomes feasible.
For the Condorcet requirement, the cost is that of a sufficient condition and therefore an upper bound for any particular Condorcet-consistent rule (Schulze, Kemeny, \dots); it is not that rule's exact cost.

\paragraph{Caps, scales and cohort (modelling assumptions).}
The cap is the central assumption and the headline result is investment as a function of it.
We use a loose ceiling ($c_t=100$), a frontier cap (best observed score on $t$ plus a margin, never below the current score), and an empirical cap (the current score plus a high percentile of observed fine-tune-minus-base score gains on $t$, never below the current score and never above the ceiling); further caps appear in the cap sweep.
Frontier and empirical caps are assumptions, not evidence that such scores are attainable; empirical caps carry survivorship bias and use one gain quantile for all targets.
Open LLM Leaderboard scores are on unrelated scales, so raw scores, per-task standardised scores and rank percentiles are distinct variants (Section~\ref{sec:exp}).
Geometric means use only positive-domain scales; $z$-scores are used only for mean, median, minimum and pooled $L_1$.
Transform parameters are fixed from the original snapshot and edits and caps are mapped through the same transform.
In every comparison across rules the targets, rivals and caps are identical (a fixed cohort); where the ordinal solvers are too slow, a rank-stratified sample is used for \emph{all} rules in that table and results are reported by rank bin.
Task bootstraps and evaluation-noise simulations are sensitivity analyses conditional on these assumptions and do not validate them.

\paragraph{Closed forms for the score rules.}
Let $\tau$ be the best rival's value under the rule, $b=\tau+\varepsilon$, $h_t=\max(c_t-x_{it},0)$ the headroom, and $x=x_{i\cdot}$.
Write $d^\ast$ for the smallest $d$ such that $y=\min(x+d,c)$ makes the target first: since every rule is monotone in the target's own scores, any set of edits that wins can be replaced by this uniform-lift edit with no larger largest lift, so $L_\infty^\ast=d^\ast$ for every rule (Proposition~4).
\begin{enumerate}\setlength\itemsep{0pt}
\item \textbf{Mean.} Let $\mathrm{need}=nb-\sum_t x_t$. The target is already first if $\mathrm{need}\le 0$; it is feasible iff $\sum_t h_t\ge\mathrm{need}$; then $L_1^\ast=\mathrm{need}$ and $K^\ast$ is the fewest tasks whose largest headrooms sum to at least $\mathrm{need}$.
\item \textbf{Median, $n$ odd.} The target needs at least $m=(n+1)/2$ tasks with $y_t\ge b$. With $a=\#\{t:x_t\ge b\}$, $K^\ast=\max(0,m-a)$; $L_1^\ast$ is the sum of the $K^\ast$ smallest values of $b-x_t$ over tasks below $b$ with $c_t\ge b$; it is infeasible iff fewer than $K^\ast$ such tasks exist. (For even $n$, as on the six-task Open LLM Leaderboard, the median is the mean of two middle values and the solver is used; it is checked against brute force.)
\item \textbf{Minimum.} $K^\ast=\#\{t:x_t<b\}$, $L_1^\ast=\sum_t(b-x_t)^+$, feasible iff $c_t\ge b$ on every task below $b$.
\end{enumerate}
These formulas explain the cap dependence we find: a rule that requires every task to reach $b$ (minimum) or a majority of tasks (median) is blocked as soon as the cap lies below $b$ on those tasks, whereas the mean can trade headroom across tasks.
They are conditional on the model above; they are not statements about training.
We verify the formulas, and Proposition~4 for the geometric mean, Borda and Condorcet as well, against the exact solvers on task subsets, targets and caps sampled from the two matrices (31{,}080 checks, all agreeing; Appendix, Table~\ref{tab:theory}); the proofs are short and are given there.

\section{Experiments}
\label{sec:exp}

\paragraph{Data and cohorts.}
OLLB: the \texttt{contents} table of the Open LLM Leaderboard v2 (pinned commit \texttt{9c09a7ca}), with its six official normalised task scores (IFEval, BBH, MATH Lvl~5, GPQA, MuSR, MMLU-Pro), 4{,}496 models after removing 79 duplicate names and one flagged model (submissions 2024-06 to 2025-03).
HELM MMLU: release v1.13.0, 79 models, exact-match accuracy on 57 subjects.
Both are frozen snapshots; all ranks and costs refer to them.
A \emph{target} is any model of rank $\ge2$ under the mean rule: 4{,}495 on OLLB and 78 on HELM.
Score rules run on every target.
The ordinal solvers are slower, so Borda and Condorcet run on a rank-stratified sample of 136 OLLB targets (ranks 2--21 and log-uniform tail draws) with stratum weights that make weighted shares and medians estimate the full-leaderboard quantity; HELM runs every rule on every target.
We report by rank group of the target under the mean rule (2--10, 11--100, $>$100).
Caps (Section~\ref{sec:method}): ceiling (100), frontier${}+2$ (best rival${}+2$ points), $+5/10/20$ points over the target's own score, and the empirical 95th and 99th percentiles of fine-tune-minus-base gains (OLLB, 540 base/fine-tune pairs).
For HELM there are no such pairs, so empirical caps use the OLLB MMLU-Pro gain percentile for every subject, a proxy marked $^\dagger$ throughout.
Intervals are 95\% percentile bootstraps over targets within rank strata; they describe these snapshots, not a population of leaderboards.
OLLB values are normalised points and HELM values percentage points, so costs are compared across rules, never across datasets.
The solvers are exact (closed forms, convex or greedy solvers, and HiGHS integer programs), and every solver status is either recorded as optimal or cross-checked against an independent feasibility test (Appendix~\ref{app:repro}).

\paragraph{Headline: how many targets can reach first place at all.}
Figure~\ref{fig:feas} shows the share of targets for which first place is attainable when no task may be raised by more than $h$ points, i.e.\ the distribution of the smallest uniform lift $d^\ast$ (a target is feasible under a gain cap $h$ exactly when $d^\ast\le h$, which we verified on every row).
Table~\ref{tab:feas} adds the empirical cap and the rank groups.
\begin{figure}[!tbp]
\centering
\includegraphics[width=0.8\linewidth]{../figures/F1_feasibility_curve.pdf}
\caption{Share of targets for which first place is attainable when no task may be raised by more than $h$ points (weighted empirical CDF of the minimal uniform lift $d^\ast$; vertical lines at $h=5,10,20$). Top: all targets of rank $\ge2$ (OLLB 4{,}495; HELM 78); bottom: the 29 targets ranked 2--30. Borda and Condorcet on OLLB: 136-target stratified sample, stratum-weighted.}
\label{fig:feas}
\end{figure}
\begin{table}[!tbp]
\centering\footnotesize
\caption{Feasibility of first place by rule and the target's rank under the mean (2--10, 11--100, $>$100): share of targets for which first place is attainable with at most 10 points per task, and under the empirical 95th-percentile gain cap; and the minimal uniform lift $d^\ast$ in points, median [95\% bootstrap interval, omitted where narrower than 1 point]. OLLB: 4{,}495 targets (score rules), 136-target stratified sample (Borda, Condorcet); HELM: 78 targets, none ranked beyond 79. $^\dagger$HELM empirical caps are a proxy.}
\label{tab:feas}
\begingroup\input{../figures/tables/T1_feasibility.tex}\endgroup
\end{table}
At a 10-point cap the shares of all OLLB targets are 3\% (mean), 5\% (median), 10\% (minimum), 3\% (geometric mean), 29\% (Borda) and 3\% (Condorcet winner); on HELM they are 37\%, 51\%, 18\%, 31\%, 31\% and 49\%.
Almost all models far from the top cannot reach first place: of the OLLB targets ranked beyond 100 at most 28\% (Borda; 9\% for the minimum) and under the empirical 95\% cap at most 25\% can (Table~\ref{tab:feas}).
For the 29 best-placed targets (ranks 2--30) the shares are high, 76--100\% on OLLB and 41--100\% on HELM, and the minimum rule is lowest on both.
The rules change \emph{who} can win, not only the price: at a 20-point cap the minimum, Borda and Condorcet rules are attainable for 46--64\% of OLLB targets against 14--21\% for the mean, median and geometric mean, whereas on HELM the minimum rule is the hardest at 5, 10 and 20 points (9\%, 18\% and 44\%).
The empirical cap is stricter: under the 95th-percentile cap only the targets near the top qualify (OLLB ranks 2--10: 100\% for mean, median, geometric mean and Borda, 44--56\% for minimum and Condorcet; ranks 11--100: 22--100\%), and on HELM, with the proxy cap, at most 1\% of targets of rank 11--100 for any rule.
The smallest uniform lift grows with rank under every rule (OLLB median $d^\ast$ of 0--8 points for ranks 2--10, 5--9 for ranks 11--100, 18--36 beyond rank 100; HELM 11--23 points for ranks 11--100).

\begin{figure}[!tbp]
\centering
\includegraphics[width=0.74\linewidth]{../figures/F2_cap_sweep.pdf}
\caption{Cap sweep. Top: share of targets (rank $\ge2$) for which first place is attainable, by rule and cap scenario (caps are not nested; $^\dagger$ proxy). Bottom: median $K^\ast$ on the common feasible cohort of that cap (tick label: cohort size / targets run); points are shown for cohorts of at least 5 targets, bars: 95\% bootstrap intervals. Cohorts differ between caps: compare rules within one cap.}
\label{fig:cap}
\end{figure}

\paragraph{Cost on the common feasible cohort.}
Costs are compared only among targets feasible under every listed rule at that cap (Table~\ref{tab:cost}).
\begin{table}[!tbp]
\centering\footnotesize
\caption{Cost of first place on the common feasible cohort at a cap, median [95\% bootstrap interval] of tasks touched $K^\ast$ and total points $L_1^\ast$. The cohort and the targets run are in each block header. OLLB values are normalised points, HELM values percentage points, so they are not comparable across datasets; cohorts differ between blocks, so compare rules within a block. The ceiling and frontier${}+2$ caps leave every target feasible; cohorts at attainable caps are small and are in Appendix~\ref{app:tables}.}
\label{tab:cost}
\begingroup\input{../figures/tables/T2_common_cohort_cost.tex}\endgroup
\end{table}
At the ceiling cap on OLLB the minimum rule needs the fewest points (25.3 against 100.6 for the mean), then the median (48.1), Borda (65.6), Condorcet (78.1) and the geometric mean (78.8).
On HELM the minimum (98) and median (111) rules need the fewest points and the mean (661) and geometric mean (633) the most, with Condorcet (213) and Borda (500) in between; relative to the mean's $L_1^\ast$, the median rule's median is 0.15 (95\% interval 0.14--0.19) and the minimum's 0.14 on HELM, and 0.60 and 0.25 on OLLB (\texttt{rel\_L1\_vs\_mean}).
The number of tasks tells a different story: Borda needs the most tasks (OLLB 5 of 6; HELM 46 of 57), the minimum rule the fewest on HELM (11), and the mean rule's $K^\ast$ depends strongly on the cap, because a binding cap forces the total to be spread over more tasks (HELM: 16 at the ceiling, 31 at frontier${}+2$).
So the rule that is cheapest in points is not the one that touches the fewest tasks, and the ranking of rules depends on which investment measure is used; on both datasets the ordinal rules need many tasks but not many points relative to the mean.

\paragraph{Sensitivity to the cap, rank, number of tasks and scale.}
Figure~\ref{fig:cap} sweeps all cap scenarios.
Ceiling and frontier caps leave every target feasible, while at $+5$ points 0--2\% of OLLB targets are feasible and at $+20$ points 14--64\%, so the feasibility ordering of the rules shown above is what changes with the cap; we do not rank costs across rules at attainable caps, where cohorts are small (3 to 91 targets) and differ.
Cost rises with the target's rank for every rule overall, though not strictly at every rank (Appendix Figure~\ref{fig:rank}).
On HELM random subsets of $n$ subjects (100 draws per $n$) show that Borda and Condorcet need more tasks as $n$ grows, so part of the Borda cost above reflects $n=57$ (Appendix Figure~\ref{fig:subset}).
On OLLB, replacing the official normalised scores with raw accuracies, per-task $z$-scores or rank percentiles changes median $K^\ast$ for targets ranked 3--10 by at most one task, but for targets ranked beyond 10 by up to four: at the ceiling the mean rule needs 2 tasks on the official scale and 6 on rank percentiles, and the minimum rule 3 against 5 (raw) and 6 ($z$-score) (Appendix Figure~\ref{fig:scale}). We therefore state OLLB conclusions for the official scale and treat the others as sensitivity.


\paragraph{Robustness and theory check.}
Changing the strictness margin $\varepsilon$ from $10^{-3}$ to $10^{-2}$ or $10^{-1}$ changes the median $L_1^\ast$ by at most 12\%; $K^\ast$ changes for at most 13\% of targets at the ceiling cap (Table~\ref{tab:eps}).
Removing merged models leaves median $K^\ast$ at the ceiling and frontier caps unchanged; keeping only the 500 top-ranked models leaves it unchanged for targets ranked 3--10 and lowers it by one to two tasks for lower-ranked targets (Table~\ref{tab:slices}).
An evaluation-noise Monte Carlo and a task bootstrap on 10 targets (100--200 draws) move the feasible share of the minimum rule from the observed 70\% to 49--70\% and leave median $K^\ast$ within one task of the observed value on OLLB and within up to eight tasks (Borda) on HELM (Table~\ref{tab:noise}); these are conditional sensitivity analyses on few targets, not estimates of leaderboard sampling variation.
The closed forms (mean, median for odd $n$, minimum) and the uniform-lift identity agree with the solvers on all 31{,}080 instances (Table~\ref{tab:theory}).

\section{Discussion and Limitations}
\label{sec:limits}

\paragraph{What the results say.}
Under attainable caps, most models cannot reach first place under any of the six rules, and the rule determines how many can: the share at 10 points ranges from 3\% to 29\% on OLLB and from 18\% to 51\% on HELM.
Cost comparisons across rules are meaningful only where the same targets are feasible under all of them; there, the minimum and median rules need the fewest points, Borda the most tasks, and the mean rule's task count depends on the cap.
The ordering of rules differs between OLLB and HELM and between caps (Figures~\ref{fig:feas} and~\ref{fig:cap}), so the supported statement is conditional: it holds per cap scenario and scale, not as a single ranking of rules.
The closed forms identify why: the minimum rule needs every task below the best rival's level to reach it, the median rule a majority, and the mean rule only a total, which any cap that does not bind can supply from few tasks.

\paragraph{Limitations.}
\emph{The estimand is a score counterfactual, not a training cost.}
It assumes that any score in $[x_{it},c_t]$ is reachable independently on each task, with no spillover (gains on one task need not transfer, and may reduce other scores), no rival response and no equilibrium, and it says nothing about the data, compute or contamination that would deliver the gain.
\emph{Caps are assumptions.}
Frontier and empirical caps are not evidence that such scores are attainable; the empirical caps come from observed fine-tune-minus-base pairs (survivorship bias, one quantile for all targets), and on HELM they are a proxy borrowed from another benchmark.
The conclusions on feasibility are therefore as good as the cap scenario; we show eight and rank nothing across them.
\emph{Scale and size.}
OLLB has $n=6$ tasks on unrelated scales, so median, minimum and geometric-mean results are tied to the official normalisation (Appendix Figure~\ref{fig:scale}) and the median and Condorcet results are tie-sensitive; HELM is one benchmark with 57 same-unit subjects and 79 models, and describes a snapshot, not a population of leaderboards.
\emph{Cohorts and intervals.}
Borda and Condorcet on OLLB are computed on a 136-target stratified sample (stratum-weighted); cohorts at attainable caps are small (3 to 91 targets); the Condorcet cost is an upper bound for any Condorcet-consistent rule; $K^\ast$ has few levels at $n=6$.
The task bootstrap and the evaluation-noise simulations use 10 targets and are sensitivity analyses conditional on the model and caps; they do not validate them.
\emph{Margin.}
Strict victory over a continuous score has an infimum, so we fix $\varepsilon=10^{-3}$ and report sensitivity to it (Table~\ref{tab:eps}).
\emph{Measurement and data.}
Scores are single-run leaderboard numbers; one OLLB model per name is kept; the licences of the score data are not declared on the sources we could verify.
We use no evidence about any model's training data, and nothing here accuses any listed model.

\paragraph{Broader impact.}
The study documents how little targeted improvement some rules require and how others are out of reach, which can inform more robust aggregation and also show where leaderboards are easiest to game.
It uses public score tables only, with no human subjects and no model training.

\paragraph{Conclusion.}
First place under a multi-task rule is not a single price: it is a feasibility question first and a cost question only among targets for which first place is attainable.
On the two matrices studied, the aggregation rule changes how many models can win and what it costs, and neither ordering is stable across datasets and caps.


\bibliographystyle{plainnat}
\bibliography{references}

\appendix
\section{Proofs of the Propositions}
\label{app:proofs}

Notation as in Section~\ref{sec:method}: $x=x_{i\cdot}$ is the target's score vector, $c_t\ge x_t$ the cap, $h_t=c_t-x_t\ge0$ the headroom, $\tau$ the best rival's value under the score rule, $b=\tau+\varepsilon$; the target is first iff its value is at least $b$.
Let $y$ be an edit with $x\le y\le c$.

\paragraph{Proposition 1 (mean).}
The target is first iff $\sum_t y_t\ge nb$. With $\mathrm{need}=nb-\sum_t x_t$: if $\mathrm{need}\le0$ no edit is needed. Otherwise $\sum_t(y_t-x_t)\ge\mathrm{need}$ is necessary, so $L_1^\ast\ge\mathrm{need}$; it is attainable iff $\sum_t h_t\ge\mathrm{need}$ (raise any tasks within their headroom until the total is met), which is also the feasibility condition. $K^\ast$ is the least $K$ such that the $K$ largest headrooms sum to at least $\mathrm{need}$, since an edit touching $K$ tasks adds at most the sum of their headrooms. \hfill$\square$

\paragraph{Proposition 2 (median, $n$ odd).}
The median is at least $b$ iff at least $m=(n+1)/2$ entries are $\ge b$. Let $a=\#\{t:x_t\ge b\}$; those entries already count. Every other counted task must be raised to at least $b$ and so must have $c_t\ge b$ and cost at least $b-x_t$. Hence at least $\max(0,m-a)$ tasks must be touched, $K^\ast=\max(0,m-a)$ if at least that many tasks below $b$ have $c_t\ge b$, and the target is infeasible otherwise. Raising a chosen task to exactly $b$ suffices, so $L_1^\ast$ is the sum of the $K^\ast$ smallest values of $b-x_t$ among those tasks. \hfill$\square$

\paragraph{Proposition 3 (minimum).}
The minimum is at least $b$ iff every $y_t\ge b$. Every task with $x_t<b$ must be raised, to at least $b$, which needs $c_t\ge b$; raising each to exactly $b$ suffices. So $K^\ast=\#\{t:x_t<b\}$, $L_1^\ast=\sum_t(b-x_t)^+$, and the target is feasible iff $c_t\ge b$ on every such task. \hfill$\square$

\paragraph{Proposition 4 ($L_\infty^\ast$ equals the smallest uniform lift).}
Each rule we study is monotone in the following sense: if $y$ makes the target first, so does every $y'$ with $y\le y'\le c$.
For the score rules this holds because they are nondecreasing in each coordinate and the rivals are fixed. For Borda and the Condorcet requirement, raising the target's score on a task can only improve its position against each rival on that task and cannot improve any rival's, so the target's total never falls and each rival's total never rises.
Let $y$ win with largest lift $\delta=\max_t(y_t-x_t)$ and put $y'_t=\min(x_t+\delta,c_t)$. Then $y\le y'\le c$, so $y'$ wins, and its largest lift is at most $\delta$. Hence $L_\infty^\ast=d^\ast:=\min\{d:\ \min(x+d,c)\text{ wins}\}$ for every such rule. \hfill$\square$
\emph{Corollary.} For the gain cap $c_t=\min(x_t+h,100)$, a target is feasible iff $d^\ast\le h$ (with $d^\ast$ computed at the ceiling), because by monotonicity feasibility is equivalent to winning at $y=c$.
The numerical agreement of Propositions 1--4 with the solvers is in Table~\ref{tab:theory}.

\section{Additional Tables}
\label{app:tables}

\begin{table}[h]
\centering\footnotesize
\caption{Closed forms (P1 mean, P2 median, P3 minimum; P4: $L_\infty^\ast$ equals the smallest uniform lift) against the solvers on 31{,}080 instances sampled from the two matrices (task subsets, targets, caps): all agree; largest absolute differences are at most $10^{-6}$ (tolerance $10^{-4}$); feasibility status agrees on every instance.}
\label{tab:theory}
\begingroup\input{../figures/tables/A4_theory_vs_solver.tex}\endgroup
\end{table}

\begin{table}[h]
\centering\footnotesize
\caption{Sensitivity to the strictness margin $\varepsilon$ (base $10^{-3}$): median ratio of $L_1^\ast$ for $\varepsilon=0.01$ and $0.1$ to the base, and the share of targets whose $K^\ast$ changes. ``feas.'' is the feasible share at the base margin (OLLB Borda and Condorcet: stratified sample). $^\dagger$ proxy cap.}
\label{tab:eps}
\begingroup\input{../figures/tables/A1_eps_sensitivity.tex}\endgroup
\end{table}

\begin{table}[h]
\centering\footnotesize
\caption{Cost of first place on the common feasible cohort at attainable caps (+10 points; empirical 95\%$^\dagger$), median [95\% bootstrap interval]. Cohorts are 4--25 targets; blocks with fewer than 10 targets are not summarised. $^\dagger$ proxy cap on HELM.}
\label{tab:cost_attain}
\begingroup\input{../figures/tables/A5_common_cohort_cost_attainable_caps.tex}\endgroup
\end{table}

\begin{table}[p]
\centering\scriptsize
\caption{Leaderboard slices (OLLB score rules): all models, merged models removed, top-500 models only; ranks 3--10 hold 8 targets in every slice. At +10 points, medians are among feasible targets only (feasible share beside them) and are not ranked across rules.}
\label{tab:slices}
\begingroup\input{../figures/tables/A2_slices.tex}\endgroup
\end{table}

\begin{table}[t]
\centering\footnotesize
\caption{Evaluation-noise Monte Carlo (100 draws) and task bootstrap (OLLB 200, HELM 100 draws) on 10 targets at the +10-point cap: feasible share and median $K^\ast$ [IQR] over draws against the unperturbed values for the same targets. Conditional sensitivity analyses on few targets, not validation of the cap assumption.}
\label{tab:noise}
\begingroup\input{../figures/tables/A3_noise_taskboot.tex}\endgroup
\end{table}

\section{Additional Figures}
\label{app:figs}

\begin{figure}[h]
\centering
\includegraphics[width=0.92\linewidth]{../figures/F3_cost_vs_rank.pdf}
\caption{Minimum investment against the target's rank under the mean, cap frontier${}+2$ (every target feasible). Top: $K^\ast$; bottom: $L_1^\ast$ (symlog axis; zero means already first). Lines: median over rank strata (2, 3--5, 6--10, 11--30, 31--100, $>$100), bands: 95\% bootstrap interval. Left: OLLB (Borda and Condorcet on the 136-target sample, so their strata hold few targets); right: HELM MMLU. Cost rises with rank for every rule overall, not monotonically for the median and Borda $L_1^\ast$ on OLLB at ranks 2--10, where intervals are wide.}
\label{fig:rank}
\end{figure}

\begin{figure}[h]
\centering
\includegraphics[width=0.92\linewidth]{../figures/F4_subset_size.pdf}
\caption{HELM MMLU with random subsets of $n$ subjects (100 draws per $n$): for the model ranked 2nd (left) or 5th (right) under the rule itself in the subset, median $K^\ast$ with IQR at the ceiling cap (top) and feasible share at the $+10$-point cap (bottom). One target per rule, draw and rank, so targets differ between rules. Borda and Condorcet need more tasks as $n$ grows; for the 5th-ranked model the minimum rule is the only one whose feasible share falls clearly below 100\%. Conditional sensitivity analysis over task subsets.}
\label{fig:subset}
\end{figure}

\begin{figure}[h]
\centering
\includegraphics[width=0.92\linewidth]{../figures/F5_ollb_scale_variants.pdf}
\caption{OLLB score rules on four versions of the task scores (official normalised, raw accuracy, per-task $z$-score, rank percentile): median $K^\ast$ with 95\% bootstrap interval, for targets ranked 3--10 (top; 8 targets) and $>$10 (bottom), ceiling cap (left) and frontier cap (right). All targets feasible; the geometric mean is not in the $z$-score set. Caps are in each variant's own units, so $L_1^\ast$ is not compared across scales.}
\label{fig:scale}
\end{figure}

\section{Reproducibility, Data and Use of LLMs}
\label{app:repro}

\paragraph{Reproducibility.}
The supplementary code package contains the solvers, the data build with SHA-256 checksums, all result files, the unit tests (solvers against brute force) and a reproduction guide with a script-to-table map.
Every number in this paper is read from its result files; none is typed by hand.
Environment: Python 3.11.2, numpy 2.4.6, scipy 1.17.1, pandas 3.0.6, HiGHS through \texttt{scipy.optimize.milp}; seed 0 throughout (target sampling, bootstrap, noise draws, theory instances); CPU only.
Data: OLLB \texttt{contents} at commit \texttt{9c09a7ca}\ldots (and \texttt{results} at \texttt{aa81ecc3}\ldots for the evaluation-noise model) and HELM MMLU release v1.13.0, both public and pinned.
Compute: wall-clock on a shared host, 2-CPU jobs under 1\,GB each: HELM, all rules and caps, 862\,s; OLLB score rules on all targets, 2{,}164\,s; ordinal integer programs on the 136-target sample, 1{,}698\,s (Borda) and 890\,s (Condorcet) at the frontier${}+2$ cap and 1{,}824\,s (Borda) at the ceiling; at $M=4{,}496$ a Borda solve takes a median of 8.3\,s (maximum 115\,s).
The ordinal runs at frontier${}+2$ record their solver status (all optimal); older ordinal runs did not, and for them an independent feasibility test agrees with the solver's finite/infinite answer on all of 564{,}000 checked rows (one early disagreement was re-solved), so their budgets are ``solver-reported minimum, gap not recorded''.
All 350 data-set/cap/rule cells are complete.

\paragraph{Licences and assets.}
The OLLB \texttt{contents} dataset carries no licence declaration on its card that we could verify, and HELM outputs are published by Stanford CRFM in a public bucket without a per-file licence that we checked; neither is redistributed beyond derived score matrices (model names and scores).
OLLB's per-model licence column is blank for 1{,}736 of 4{,}497 models (before removal of the one flagged model) and is not a licence for the score data.
The derived score matrices are part of the code package.

\paragraph{Use of LLMs.}
All code, experiments, analysis and text, including this manuscript, were produced by LLM agents (Claude, through Claude Code) acting as a research organisation; humans did not edit them.
We did not search for or read any publication on this exact question.


\section*{NeurIPS Paper Checklist}




\begin{enumerate}

\item {\bf Claims}
    \item[] Question: Do the main claims made in the abstract and introduction accurately reflect the paper's contributions and scope?
    \item[] Answer: \answerYes{}
    \item[] Justification: The abstract and introduction state the setting, the three investment measures, the feasibility-first findings and that conclusions are conditional on the cap and scale (Sections~\ref{sec:intro}, \ref{sec:exp}, \ref{sec:limits}); investment is declared a score counterfactual, not a training cost.
    \item[] Guidelines:
    \begin{itemize}
        \item The answer \answerNA{} means that the abstract and introduction do not include the claims made in the paper.
        \item The abstract and/or introduction should clearly state the claims made, including the contributions made in the paper and important assumptions and limitations. A \answerNo{} or \answerNA{} answer to this question will not be perceived well by the reviewers. 
        \item The claims made should match theoretical and experimental results, and reflect how much the results can be expected to generalize to other settings. 
        \item It is fine to include aspirational goals as motivation as long as it is clear that these goals are not attained by the paper. 
    \end{itemize}

\item {\bf Limitations}
    \item[] Question: Does the paper discuss the limitations of the work performed by the authors?
    \item[] Answer: \answerYes{}
    \item[] Justification: Section~\ref{sec:limits} lists the score-counterfactual assumptions (no spillover, fixed rivals), caps as assumptions, $n=6$ and unrelated OLLB scales, one HELM benchmark, snapshots, sampled cohorts, the margin $\varepsilon$ and the conditional nature of the bootstrap and noise analyses.
    \item[] Guidelines:
    \begin{itemize}
        \item The answer \answerNA{} means that the paper has no limitation while the answer \answerNo{} means that the paper has limitations, but those are not discussed in the paper. 
        \item The authors are encouraged to create a separate ``Limitations'' section in their paper.
        \item The paper should point out any strong assumptions and how robust the results are to violations of these assumptions (e.g., independence assumptions, noiseless settings, model well-specification, asymptotic approximations only holding locally). The authors should reflect on how these assumptions might be violated in practice and what the implications would be.
        \item The authors should reflect on the scope of the claims made, e.g., if the approach was only tested on a few datasets or with a few runs. In general, empirical results often depend on implicit assumptions, which should be articulated.
        \item The authors should reflect on the factors that influence the performance of the approach. For example, a facial recognition algorithm may perform poorly when image resolution is low or images are taken in low lighting. Or a speech-to-text system might not be used reliably to provide closed captions for online lectures because it fails to handle technical jargon.
        \item The authors should discuss the computational efficiency of the proposed algorithms and how they scale with dataset size.
        \item If applicable, the authors should discuss possible limitations of their approach to address problems of privacy and fairness.
        \item While the authors might fear that complete honesty about limitations might be used by reviewers as grounds for rejection, a worse outcome might be that reviewers discover limitations that aren't acknowledged in the paper. The authors should use their best judgment and recognize that individual actions in favor of transparency play an important role in developing norms that preserve the integrity of the community. Reviewers will be specifically instructed to not penalize honesty concerning limitations.
    \end{itemize}

\item {\bf Theory assumptions and proofs}
    \item[] Question: For each theoretical result, does the paper provide the full set of assumptions and a complete (and correct) proof?
    \item[] Answer: \answerYes{}
    \item[] Justification: Propositions 1--4 are stated with their assumptions in Section~\ref{sec:method} and proved in Appendix~\ref{app:proofs}; they are also checked against exact solvers on 31{,}080 instances (Table~\ref{tab:theory}).
    \item[] Guidelines:
    \begin{itemize}
        \item The answer \answerNA{} means that the paper does not include theoretical results. 
        \item All the theorems, formulas, and proofs in the paper should be numbered and cross-referenced.
        \item All assumptions should be clearly stated or referenced in the statement of any theorems.
        \item The proofs can either appear in the main paper or the supplemental material, but if they appear in the supplemental material, the authors are encouraged to provide a short proof sketch to provide intuition. 
        \item Inversely, any informal proof provided in the core of the paper should be complemented by formal proofs provided in appendix or supplemental material.
        \item Theorems and Lemmas that the proof relies upon should be properly referenced. 
    \end{itemize}

    \item {\bf Experimental result reproducibility}
    \item[] Question: Does the paper fully disclose all the information needed to reproduce the main experimental results of the paper to the extent that it affects the main claims and/or conclusions of the paper (regardless of whether the code and data are provided or not)?
    \item[] Answer: \answerYes{}
    \item[] Justification: Sections~\ref{sec:method}--\ref{sec:exp} define rules, caps, cohorts and intervals; Appendix~\ref{app:repro} gives software versions, data pins, seeds and the checked solver statuses.
    \item[] Guidelines:
    \begin{itemize}
        \item The answer \answerNA{} means that the paper does not include experiments.
        \item If the paper includes experiments, a \answerNo{} answer to this question will not be perceived well by the reviewers: Making the paper reproducible is important, regardless of whether the code and data are provided or not.
        \item If the contribution is a dataset and\slash or model, the authors should describe the steps taken to make their results reproducible or verifiable. 
        \item Depending on the contribution, reproducibility can be accomplished in various ways. For example, if the contribution is a novel architecture, describing the architecture fully might suffice, or if the contribution is a specific model and empirical evaluation, it may be necessary to either make it possible for others to replicate the model with the same dataset, or provide access to the model. In general. releasing code and data is often one good way to accomplish this, but reproducibility can also be provided via detailed instructions for how to replicate the results, access to a hosted model (e.g., in the case of a large language model), releasing of a model checkpoint, or other means that are appropriate to the research performed.
        \item While NeurIPS does not require releasing code, the conference does require all submissions to provide some reasonable avenue for reproducibility, which may depend on the nature of the contribution. For example
        \begin{enumerate}
            \item If the contribution is primarily a new algorithm, the paper should make it clear how to reproduce that algorithm.
            \item If the contribution is primarily a new model architecture, the paper should describe the architecture clearly and fully.
            \item If the contribution is a new model (e.g., a large language model), then there should either be a way to access this model for reproducing the results or a way to reproduce the model (e.g., with an open-source dataset or instructions for how to construct the dataset).
            \item We recognize that reproducibility may be tricky in some cases, in which case authors are welcome to describe the particular way they provide for reproducibility. In the case of closed-source models, it may be that access to the model is limited in some way (e.g., to registered users), but it should be possible for other researchers to have some path to reproducing or verifying the results.
        \end{enumerate}
    \end{itemize}


\item {\bf Open access to data and code}
    \item[] Question: Does the paper provide open access to the data and code, with sufficient instructions to faithfully reproduce the main experimental results, as described in supplemental material?
    \item[] Answer: \answerYes{}
    \item[] Justification: Data are public and pinned (Appendix~\ref{app:repro}); an anonymised supplementary code package holds the code, derived score matrices, results and a command-by-command guide. The source data licences are undeclared and said so.
    \item[] Guidelines:
    \begin{itemize}
        \item The answer \answerNA{} means that paper does not include experiments requiring code.
        \item Please see the NeurIPS code and data submission guidelines (\url{https://neurips.cc/public/guides/CodeSubmissionPolicy}) for more details.
        \item While we encourage the release of code and data, we understand that this might not be possible, so \answerNo{} is an acceptable answer. Papers cannot be rejected simply for not including code, unless this is central to the contribution (e.g., for a new open-source benchmark).
        \item The instructions should contain the exact command and environment needed to run to reproduce the results. See the NeurIPS code and data submission guidelines (\url{https://neurips.cc/public/guides/CodeSubmissionPolicy}) for more details.
        \item The authors should provide instructions on data access and preparation, including how to access the raw data, preprocessed data, intermediate data, and generated data, etc.
        \item The authors should provide scripts to reproduce all experimental results for the new proposed method and baselines. If only a subset of experiments are reproducible, they should state which ones are omitted from the script and why.
        \item At submission time, to preserve anonymity, the authors should release anonymized versions (if applicable).
        \item Providing as much information as possible in supplemental material (appended to the paper) is recommended, but including URLs to data and code is permitted.
    \end{itemize}


\item {\bf Experimental setting/details}
    \item[] Question: Does the paper specify all the training and test details (e.g., data splits, hyperparameters, how they were chosen, type of optimizer) necessary to understand the results?
    \item[] Answer: \answerYes{}
    \item[] Justification: There is no training; the setting is the rule definitions, tie handling, margin $\varepsilon$, cohorts and strata, caps, bootstrap and solver details in Sections~\ref{sec:method}--\ref{sec:exp} and Appendix~\ref{app:repro}.
    \item[] Guidelines:
    \begin{itemize}
        \item The answer \answerNA{} means that the paper does not include experiments.
        \item The experimental setting should be presented in the core of the paper to a level of detail that is necessary to appreciate the results and make sense of them.
        \item The full details can be provided either with the code, in appendix, or as supplemental material.
    \end{itemize}

\item {\bf Experiment statistical significance}
    \item[] Question: Does the paper report error bars suitably and correctly defined or other appropriate information about the statistical significance of the experiments?
    \item[] Answer: \answerYes{}
    \item[] Justification: Intervals are 95\% percentile bootstraps over targets within rank strata (Section~\ref{sec:exp}); the solutions are deterministic on a frozen matrix, and the task-bootstrap and noise analyses are labelled conditional sensitivity analyses, not sampling uncertainty of a leaderboard (Table~\ref{tab:noise}).
    \item[] Guidelines:
    \begin{itemize}
        \item The answer \answerNA{} means that the paper does not include experiments.
        \item The authors should answer \answerYes{} if the results are accompanied by error bars, confidence intervals, or statistical significance tests, at least for the experiments that support the main claims of the paper.
        \item The factors of variability that the error bars are capturing should be clearly stated (for example, train/test split, initialization, random drawing of some parameter, or overall run with given experimental conditions).
        \item The method for calculating the error bars should be explained (closed form formula, call to a library function, bootstrap, etc.)
        \item The assumptions made should be given (e.g., Normally distributed errors).
        \item It should be clear whether the error bar is the standard deviation or the standard error of the mean.
        \item It is OK to report 1-sigma error bars, but one should state it. The authors should preferably report a 2-sigma error bar than state that they have a 96\% CI, if the hypothesis of Normality of errors is not verified.
        \item For asymmetric distributions, the authors should be careful not to show in tables or figures symmetric error bars that would yield results that are out of range (e.g., negative error rates).
        \item If error bars are reported in tables or plots, the authors should explain in the text how they were calculated and reference the corresponding figures or tables in the text.
    \end{itemize}

\item {\bf Experiments compute resources}
    \item[] Question: For each experiment, does the paper provide sufficient information on the computer resources (type of compute workers, memory, time of execution) needed to reproduce the experiments?
    \item[] Answer: \answerYes{}
    \item[] Justification: CPU only; wall-clock times, memory and job sizes per experiment are in Appendix~\ref{app:repro}.
    \item[] Guidelines:
    \begin{itemize}
        \item The answer \answerNA{} means that the paper does not include experiments.
        \item The paper should indicate the type of compute workers CPU or GPU, internal cluster, or cloud provider, including relevant memory and storage.
        \item The paper should provide the amount of compute required for each of the individual experimental runs as well as estimate the total compute. 
        \item The paper should disclose whether the full research project required more compute than the experiments reported in the paper (e.g., preliminary or failed experiments that didn't make it into the paper). 
    \end{itemize}
    
\item {\bf Code of ethics}
    \item[] Question: Does the research conducted in the paper conform, in every respect, with the NeurIPS Code of Ethics \url{https://neurips.cc/public/EthicsGuidelines}?
    \item[] Answer: \answerYes{}
    \item[] Justification: The work uses public score tables only, with no human subjects and no model training, and accuses no listed model of anything (Section~\ref{sec:limits}).
    \item[] Guidelines:
    \begin{itemize}
        \item The answer \answerNA{} means that the authors have not reviewed the NeurIPS Code of Ethics.
        \item If the authors answer \answerNo, they should explain the special circumstances that require a deviation from the Code of Ethics.
        \item The authors should make sure to preserve anonymity (e.g., if there is a special consideration due to laws or regulations in their jurisdiction).
    \end{itemize}


\item {\bf Broader impacts}
    \item[] Question: Does the paper discuss both potential positive societal impacts and negative societal impacts of the work performed?
    \item[] Answer: \answerYes{}
    \item[] Justification: Section~\ref{sec:limits} (Broader impact) notes that the results can inform more robust aggregation and also show where leaderboards are easiest to game.
    \item[] Guidelines:
    \begin{itemize}
        \item The answer \answerNA{} means that there is no societal impact of the work performed.
        \item If the authors answer \answerNA{} or \answerNo, they should explain why their work has no societal impact or why the paper does not address societal impact.
        \item Examples of negative societal impacts include potential malicious or unintended uses (e.g., disinformation, generating fake profiles, surveillance), fairness considerations (e.g., deployment of technologies that could make decisions that unfairly impact specific groups), privacy considerations, and security considerations.
        \item The conference expects that many papers will be foundational research and not tied to particular applications, let alone deployments. However, if there is a direct path to any negative applications, the authors should point it out. For example, it is legitimate to point out that an improvement in the quality of generative models could be used to generate Deepfakes for disinformation. On the other hand, it is not needed to point out that a generic algorithm for optimizing neural networks could enable people to train models that generate Deepfakes faster.
        \item The authors should consider possible harms that could arise when the technology is being used as intended and functioning correctly, harms that could arise when the technology is being used as intended but gives incorrect results, and harms following from (intentional or unintentional) misuse of the technology.
        \item If there are negative societal impacts, the authors could also discuss possible mitigation strategies (e.g., gated release of models, providing defenses in addition to attacks, mechanisms for monitoring misuse, mechanisms to monitor how a system learns from feedback over time, improving the efficiency and accessibility of ML).
    \end{itemize}
    
\item {\bf Safeguards}
    \item[] Question: Does the paper describe safeguards that have been put in place for responsible release of data or models that have a high risk for misuse (e.g., pre-trained language models, image generators, or scraped datasets)?
    \item[] Answer: \answerNA{}
    \item[] Justification: The work releases no model and no scraped data; only derived score matrices of public leaderboards.
    \item[] Guidelines:
    \begin{itemize}
        \item The answer \answerNA{} means that the paper poses no such risks.
        \item Released models that have a high risk for misuse or dual-use should be released with necessary safeguards to allow for controlled use of the model, for example by requiring that users adhere to usage guidelines or restrictions to access the model or implementing safety filters. 
        \item Datasets that have been scraped from the Internet could pose safety risks. The authors should describe how they avoided releasing unsafe images.
        \item We recognize that providing effective safeguards is challenging, and many papers do not require this, but we encourage authors to take this into account and make a best faith effort.
    \end{itemize}

\item {\bf Licenses for existing assets}
    \item[] Question: Are the creators or original owners of assets (e.g., code, data, models), used in the paper, properly credited and are the license and terms of use explicitly mentioned and properly respected?
    \item[] Answer: \answerYes{}
    \item[] Justification: The OLLB and HELM sources, versions and pins are credited in Section~\ref{sec:exp} and Appendix~\ref{app:repro}; the score-data licences are not declared on the sources we could verify and we state that instead of asserting one.
    \item[] Guidelines:
    \begin{itemize}
        \item The answer \answerNA{} means that the paper does not use existing assets.
        \item The authors should cite the original paper that produced the code package or dataset.
        \item The authors should state which version of the asset is used and, if possible, include a URL.
        \item The name of the license (e.g., CC-BY 4.0) should be included for each asset.
        \item For scraped data from a particular source (e.g., website), the copyright and terms of service of that source should be provided.
        \item If assets are released, the license, copyright information, and terms of use in the package should be provided. For popular datasets, \url{paperswithcode.com/datasets} has curated licenses for some datasets. Their licensing guide can help determine the license of a dataset.
        \item For existing datasets that are re-packaged, both the original license and the license of the derived asset (if it has changed) should be provided.
        \item If this information is not available online, the authors are encouraged to reach out to the asset's creators.
    \end{itemize}

\item {\bf New assets}
    \item[] Question: Are new assets introduced in the paper well documented and is the documentation provided alongside the assets?
    \item[] Answer: \answerYes{}
    \item[] Justification: The derived score matrices and result files are documented in the reproduction guide of the supplementary code package (Appendix~\ref{app:repro}).
    \item[] Guidelines:
    \begin{itemize}
        \item The answer \answerNA{} means that the paper does not release new assets.
        \item Researchers should communicate the details of the dataset\slash code\slash model as part of their submissions via structured templates. This includes details about training, license, limitations, etc. 
        \item The paper should discuss whether and how consent was obtained from people whose asset is used.
        \item At submission time, remember to anonymize your assets (if applicable). You can either create an anonymized URL or include an anonymized zip file.
    \end{itemize}

\item {\bf Crowdsourcing and research with human subjects}
    \item[] Question: For crowdsourcing experiments and research with human subjects, does the paper include the full text of instructions given to participants and screenshots, if applicable, as well as details about compensation (if any)? 
    \item[] Answer: \answerNA{}
    \item[] Justification: The work involves no crowdsourcing or human subjects.
    \item[] Guidelines:
    \begin{itemize}
        \item The answer \answerNA{} means that the paper does not involve crowdsourcing nor research with human subjects.
        \item Including this information in the supplemental material is fine, but if the main contribution of the paper involves human subjects, then as much detail as possible should be included in the main paper. 
        \item According to the NeurIPS Code of Ethics, workers involved in data collection, curation, or other labor should be paid at least the minimum wage in the country of the data collector. 
    \end{itemize}

\item {\bf Institutional review board (IRB) approvals or equivalent for research with human subjects}
    \item[] Question: Does the paper describe potential risks incurred by study participants, whether such risks were disclosed to the subjects, and whether Institutional Review Board (IRB) approvals (or an equivalent approval/review based on the requirements of your country or institution) were obtained?
    \item[] Answer: \answerNA{}
    \item[] Justification: The work involves no human subjects.
    \item[] Guidelines:
    \begin{itemize}
        \item The answer \answerNA{} means that the paper does not involve crowdsourcing nor research with human subjects.
        \item Depending on the country in which research is conducted, IRB approval (or equivalent) may be required for any human subjects research. If you obtained IRB approval, you should clearly state this in the paper. 
        \item We recognize that the procedures for this may vary significantly between institutions and locations, and we expect authors to adhere to the NeurIPS Code of Ethics and the guidelines for their institution. 
        \item For initial submissions, do not include any information that would break anonymity (if applicable), such as the institution conducting the review.
    \end{itemize}

\item {\bf Declaration of LLM usage}
    \item[] Question: Does the paper describe the usage of LLMs if it is an important, original, or non-standard component of the core methods in this research? Note that if the LLM is used only for writing, editing, or formatting purposes and does \emph{not} impact the core methodology, scientific rigor, or originality of the research, declaration is not required.
    %this research? 
    \item[] Answer: \answerYes{}
    \item[] Justification: LLM agents produced the code, experiments, analysis and this manuscript; this is declared in Appendix~\ref{app:repro}.
    \item[] Guidelines:
    \begin{itemize}
        \item The answer \answerNA{} means that the core method development in this research does not involve LLMs as any important, original, or non-standard components.
        \item Please refer to our LLM policy in the NeurIPS handbook for what should or should not be described.
    \end{itemize}

\end{enumerate}
\end{document}
